\begin{slikaDesno}{fig/kalibracija_B.pdf}
\PID
Сензор за мерење магнетске инудкције калибрише се помоћу танаког референтног намотаја израђеног од $N = 100$ 
кружних завојака танке бакарне жице. Полупречник намојатаја је $a = 40\unit{mm}$ а сензор, физички малих димензија, постављен 
је у центру тог прстена. Сензор мери интензитет магнетске инукције $|\bf{B}|$, а његов излаз је напонски, сразмеран том интензитету
према непонзатом коефицијенту $V_{\rm I} = k |\bf{B}|$. 
(а) Одредити непознати калибрациони коефицијент $k$, помоћу методе најмањих квадрата.
(б) Проценити апсолутну грешку тог калибрационог коефицијента, $\Delta k$.
\\
\end{slikaDesno}
\begin{table}[ht!]
    \centering
    \begin{tabular}{c||ccccccccc} 
        \hline
        $I\unit{mA}$ & 10 & 20 & 30 & 40 & 50 & 60 & 70 & 80 & 90 \\
        $V_{\rm I}\unit{mV}$ & 2,40 & 5,69 & 15,01 & 18,04 & 17,38 & 26,72 & 31,39 & 35,57 & 35,00 \\
        \hline
    \end{tabular}
    \caption{Калибрациона мерења.}
\end{table}

\RESENJE
Магнетска индукција у пољу 